\documentclass[11pt]{article}
\usepackage[margin=0.82in]{geometry}
\usepackage{amsmath,amssymb,mathtools}
\usepackage{booktabs,tabularx}
\usepackage[dvipsnames]{xcolor}
\usepackage[colorlinks=true,linkcolor=NavyBlue,urlcolor=BrickRed,hypertexnames=false]{hyperref}
\usepackage{microtype}
\usepackage{fancyhdr}
\usepackage{enumitem}

\definecolor{ink}{HTML}{25213A}
\definecolor{accent}{HTML}{7A3E65}
\definecolor{pale}{HTML}{F8EFF5}
\definecolor{passgreen}{HTML}{16794A}
\color{ink}
\pagestyle{fancy}
\fancyhf{}
\lhead{$A_5$ representations}
\rhead{Proof and verification}
\cfoot{\thepage}
\setlength{\headheight}{14pt}
\newcommand{\Sfun}{\mathcal S}
\newcommand{\Sym}{\operatorname{Sym}}
\newcommand{\PASS}{\textcolor{passgreen}{\textbf{PASS}}}
\newcommand{\summarybox}[1]{\begin{center}\fcolorbox{accent}{pale}{\parbox{0.91\linewidth}{#1}}\end{center}}

\title{\textbf{Three Representations Associated with $A_5$}\\[3pt]
\large Class spectra, invariant multiplicities, and direct matrix tests}
\author{Computational verification note}
\date{17 July 2026}

\begin{document}
\maketitle

\summarybox{\textbf{Outcome.} All 45 multivariate coefficients indexed by
partitions through total degree seven agree in each of the five-dimensional
permutation, four-dimensional deleted permutation, and three-dimensional
rotation representations.  The explicit groups all have order 60.  Separate
Reynolds projectors confirm the low-degree invariant dimensions.}

\section{The identity being tested}

For a finite matrix group $G\leq GL(V)$, put
\begin{equation}
\Sfun_{G,V}(\mathbf t)=\frac1{|G|}\sum_{g\in G}
\prod_{a\geq1}\det(I-t_ag)^{-1}.                            \tag{1}
\end{equation}
If $g$ has eigenvalues $\lambda_1,\ldots,\lambda_n$, then
\[
\det(I-tg)^{-1}=\prod_i(1-t\lambda_i)^{-1}
=\sum_{d\geq0}h_d(\lambda_1,\ldots,\lambda_n)t^d,
\]
and $h_d(\lambda(g))=\chi_{\Sym^dV}(g)$.  Hence
\begin{align}
[t_1^{\tau_1}\cdots t_\ell^{\tau_\ell}]\Sfun_{G,V}
&=\frac1{|G|}\sum_{g\in G}\prod_a\chi_{\Sym^{\tau_a}V}(g)\notag\\
&=\dim\left(\bigotimes_a\Sym^{\tau_a}V\right)^G.          \tag{2}
\end{align}
The last equality follows because the group average on the tensor product is
the Reynolds projector.  Equations (1)--(2) also prove that grouping elements
with the same spectrum is exact in every degree.

\section{Three class-spectrum formulas}

Write $F_{(\lambda_1,\ldots,\lambda_n)}(\mathbf t)
=\prod_{a,i}(1-\lambda_i t_a)^{-1}$,
$\omega=e^{2\pi i/3}$, and $\xi=e^{2\pi i/5}$.  The formulas tested are
\begin{align}
\Sfun_{\rm perm}=\frac1{60}\big[&F_{(1,1,1,1,1)}
+15F_{(1,1,1,-1,-1)}+20F_{(1,1,1,\omega,\omega^{-1})}\notag\\
&+24F_{(1,\xi,\xi^2,\xi^3,\xi^4)}\big],                  \tag{3}\\
\Sfun_{\rm del}=\frac1{60}\big[&F_{(1,1,1,1)}
+15F_{(1,1,-1,-1)}+20F_{(1,1,\omega,\omega^{-1})}\notag\\
&+24F_{(\xi,\xi^2,\xi^3,\xi^4)}\big],                   \tag{4}\\
\Sfun_{\rm rot}=\frac1{60}\big[&F_{(1,1,1)}
+15F_{(1,-1,-1)}+20F_{(1,\omega,\omega^{-1})}\notag\\
&+12F_{(1,\xi,\xi^{-1})}+12F_{(1,\xi^2,\xi^{-2})}\big]. \tag{5}
\end{align}

For (3), a permutation of cycle type $(1^5)$, $(2^2,1)$, $(3,1^2)$,
or $(5)$ has precisely the displayed roots of unity as eigenvalues.  Removing
the always-fixed all-ones line deletes one eigenvalue $1$ and proves (4).
For (5), a rotation through angle $\theta$ has spectrum
$(1,e^{i\theta},e^{-i\theta})$; the rotation counts are
$1,15,20,12,12$ for angles $0,\pi,2\pi/3,2\pi/5,4\pi/5$.

The split of the fifth turns is essential in (5): the two classes have
different 3D spectra.  They may be combined in (3)--(4) because a 5-cycle and
its nontrivial powers have the same full permutation spectrum.

\section{Independent matrix construction}

For the first representation, all 60 even permutations of five symbols are
converted directly to $5\times5$ permutation matrices.  For the deleted
representation, let $B$ be an orthonormal basis for
\[
\{x\in\mathbb C^5:x_1+\cdots+x_5=0\};
\]
the matrices are $B^*P_\pi B$.  Neither route uses the spectral tables above.

For the rotation representation, start from the twelve icosahedron vertices
\[
(0,\pm1,\pm\varphi),\quad(\pm1,\pm\varphi,0),\quad
(\pm\varphi,0,\pm1),\qquad \varphi=(1+\sqrt5)/2.
\]
Mapping a fixed independent vertex frame to every compatible target frame and
retaining the orientation-preserving isometries reconstructs 60 distinct
$3\times3$ matrices.

\section{Verification protocol}

For every partition $\tau$ through total degree seven, the direct route
computes $h_{\tau_a}$ recursively from the eigenvalues of every explicit
matrix and averages $\prod_a h_{\tau_a}$.  The formula route uses only the
weighted spectra in (3), (4), or (5).  It never receives the matrices.

As a third route, the code explicitly constructs each $\Sym^2(V)$ inside
$V^{\otimes2}$ using normalized occupancy states, averages the resulting
matrices, and checks projector rank, trace, and idempotence.  The same check is
also performed on $V\otimes V$.

\section{Results}

\begin{center}
\begin{tabularx}{\linewidth}{@{}Xrrrrc@{}}
\toprule
Representation & $|G|$ & $\dim V$ & degree & checks & result\\
\midrule
$A_5$ permutation & 60 & 5 & $0$--$7$ & 45 & \PASS\\
$A_5$ deleted permutation & 60 & 4 & $0$--$7$ & 45 & \PASS\\
$A_5$ rotation & 60 & 3 & $0$--$7$ & 45 & \PASS\\
\bottomrule
\end{tabularx}
\end{center}

\begin{center}
\begin{tabular}{@{}lrrrr@{}}
\toprule
Representation & $[m_{(1)}]$ & $[m_{(2)}]$ & $[m_{(1,1)}]$ & $[m_{(1,1,1)}]$\\
\midrule
permutation & 1 & 2 & 2 & 5\\
deleted permutation & 0 & 1 & 1 & 1\\
rotation & 0 & 1 & 1 & 1\\
\bottomrule
\end{tabular}
\end{center}

The symmetric-square projector ranks are respectively $2,1,1$.  Their traces
agree with those integers to $3\times10^{-16}$, and the largest idempotence
error is $3.27\times10^{-16}$.  These checks recover the claims that the
permutation representation has a fixed line while the deleted and rotation
representations do not.

The value $[m_{(1,1,1)}]=1$ in the rotation case is the invariant orientation
tensor in $V^{\otimes3}$.  This is not contradicted by the absence of an
invariant cubic in $\Sym^3(V)$: $(1,1,1)$ denotes three separate degree-one
factors, not the partition $(3)$.

\section{Reproducibility and scope}

Run the command-line suite from the repository root:
{\scriptsize
\begin{verbatim}
.venv/bin/python -m lambda_distributions.proofs.matrix_group_formula_verification.a5_representations.verification
\end{verbatim}
}
Open \texttt{marimo\_notebooks/a5\_representations.py} with marimo to change the
degree bound and inspect every row.

The finite computation is a transcription and convention test, not the source
of the all-degree theorem.  All-degree validity follows from (1)--(2) and the
complete class spectra; the numerical suite checks the matrix constructions,
class sizes, fifth-root choices, and coefficient extraction in representative
small dimensions.

\end{document}
